\subsection{Permutations avoiding the specified patterns are CDG} \label{subsect:backward} The remainder of this section is devoted to the backward direction of Conjecture \ref{conj:mainresult}.  We will build up to a use of \cite[Corollary 4.13]{KR}.  We begin with some notation.

If $I_w$ is the Schubert determinantal ideal of the permutation $w \in S_n$, we will use $Z_w$ to denote the matrix obtained from an $n \times n$ matrix of indeterminates by setting $z_{i,j}$ to $0$ whenever $(i,j) \in \Dom(w)$.   If $(i,j)$ is a lower outside corner of $D_w$ and $y = z_{i,j}$, we write the CDG generators of $I_w$ as $\{yq_1+r_1, \ldots, yq_k+r_k, h_1, \ldots, h_\ell\}$ where $y$ does not divide any term of any $q_i$, $r_i$ or $h_j$.  Define $N_{y,I_w} = (h_1, \ldots, h_\ell)$ and $C_{y,I_w} = (q_1, \ldots, q_k, h_1, \ldots, h_\ell)$.  This notation mimics that in \cite{KR}.  When the CDG generators are a Gr\"obner basis of $I_w$, $C_{y,I_w}$ will be the ideal corresponding to the star and $N_{y,I_w}+(y)$ the ideal corresponding to the deletion in a geometric vertex decomposition in the sense of \cite{KMY09}.

We will call $\{q_1, \ldots, q_k, h_1, \ldots, h_\ell\}$ the \emph{CDG generators} of $C_{y,I_w}$, which is itself not typically a Schubert determinantal ideal.  With notation as above, we begin by showing that $N_{y,I_w}$ is the Schubert determinantal ideal of a permutation whose Coxeter length is smaller than that of $w$, which will be an essential component of an inductive argument.

Let $t_{a,b}$ denote the transposition $(ab) \in S_n$.

\begin{lemma}\label{lem:SchubertDel}
Suppose that $I_w$ is the Schubert determinantal ideal of the permutation $w \in S_n$ and that $(i,j)$ is a lower outside corner of $D_w$ corresponding to the variable $y = z_{i,j}$.  The ideal $N_{y,I_w}$ is the Schubert determinantal ideal of a permutation $w' \in S_n$ whose Coxeter length is strictly smaller than that of $w$.  Specifically, $w' = wt_{i,w^{-1}(j)}$, and $D_w = D_{w'} \sqcup \{(i,j)\}$.  
\end{lemma}
\begin{proof}
We claim that whenever there is some $(\rank_w(i,j)+1)$-minor with $yq+r = r$, that $r \in N_{y,I_w}$, i.e. that all of the CDG generators of $I_w$ determined only by the rank condition at $(i,j)$ involve $y$.  Fix a $(\rank_w(i,j)+1) \times (\rank_w(i,j)+1)$ submatrix $M$ of $Z_w$ so that $\det(M) = yq+r$, and suppose that $q = 0$.  Recall that the entry in row $a$ and column $b$ of $M$ is $0$ if and only if $(a,b) \in \Dom(w)$.  Then $q$ is the determinant of a $\rank_w(i,j) \times \rank_w(i,j)$ submatrix of $M$ whose anti-diagonal has an entry in row $a$ and column $b$ for some $(a,b) \in \Dom(w)$.  Assume without loss of generality that $i \geq j$.  

If $r \neq 0$, then there is some $(i',j)$ with nonvanishing $z_{i',j} \cdot q'$  summand of $r$ so that $q'$ is the determinant of a $\rank_w(i,j) \times \rank_w(i,j)$ submatrix of $M$ without a $0$ along its anti-diagonal.  Because $\Dom(w)$ forms a partition shape, if $z_{i',j}\cdot q' \neq 0$, then $z_{i'',j} \cdot q'' \neq 0$ whenever $i''<i'$ and $q''$ is the cofactor corresponding to $z_{i'',j}$ in an expansion of $yq+r$ along column $j$.  In particular, there is a unique $t$ so that no $\rank_w(i,j) \times \rank_w(i,j)$ submatrix of $M$ that excludes column $j$ and involves the final $t$ rows has a $0$ along its anti-diagonal and every $\rank_w(i,j) \times \rank_w(i,j)$ submatrix of $M$ that excludes column $j$ and also excludes one of the final $t$ rows has a $0$ along its anti-diagonal.  

The same argument can be applied to columns, and, because vanishing is determined by $0$'s along the anti-diagonal, will select the final $\rank_w(i.j)+1-t$ columns.  Hence, we may write  $r$ as the product of one $t$-minor determined by the final $t$ rows and initial $t$ columns of $M$ and one $(\rank_w(i.j)+1-t)$-minor consisting of the initial $\rank_w(i.j)+1-t$ rows and final $\rank_w(i.j)+1-t$ columns of $M$.  

Call the southeast corner of the lower block $z$ and the southeast corner of the upper block $z'$.  If there are fewer than $t$ dots northwest of $z$, then the $t$-minor that is one factor of $r$ is an element of $N_{y,I_w}$, and so $r \in  N_{y,I_w}$.  If there are $t$ or more dots northwest of $z$, then there are at most $\rank_w(i.j)-t$ dots northwest of $z'$, and so the factor of $r$ corresponding to the upper block is an element of $N_{y,I_w}$, and so $r \in N_{y,I_w}$.  Hence, $N_{y,I_w}$ is generated by the CDG generators of $I_w$ determined by the diagram boxes other than $(i,j)$.  The permutation with exactly those diagram boxes and rank conditions at those diagram boxes is $w' = wt_{i,w^{-1}(j)}$.  Then $N_{y,I_w} = I_{w'}$, the Coxeter length of $w'$ is one less than the Coxeter length of $w$, and $D_w = D_{w'} \sqcup \{(i,j)\}$.  
\end{proof}
